\section*{अध्याय \ref{ch:g9:ions} --- आयन और आयनिक विलयन}

\begin{solution}{exo:g9:ions:1}
ऐसा \omterm{def:g7:atoms-and-molecules:atom}{परमाणु} या \omterm{def:g7:atoms-and-molecules:atom}{परमाणुओं} का समूह जिसने \omterm{def:g9:inside-the-atom:nucleus}{इलेक्ट्रॉन} खोए या लिए हैं और जिस पर
आवेश है। \omterm{def:g9:ions:ion}{धनायन} धनात्मक होता है (\omterm{def:g9:inside-the-atom:nucleus}{इलेक्ट्रॉन} खोए), \omterm{def:g9:ions:ion}{ऋणायन} ऋणात्मक
(\omterm{def:g9:inside-the-atom:nucleus}{इलेक्ट्रॉन} लिए)।
\end{solution}

\begin{solution}{exo:g9:ions:2}
\ce{Mg^{2+}}: 12 \omterm{def:g9:inside-the-atom:proton}{प्रोटॉन}, 10 \omterm{def:g9:inside-the-atom:nucleus}{इलेक्ट्रॉन}।
\end{solution}

\begin{solution}{exo:g9:ions:3}
\ce{KCl}, \ce{MgCl2}, \ce{CaO}।
\end{solution}

\begin{solution}{exo:g9:ions:4}
नमकीन पानी में \omterm{def:g9:ions:ion}{आयन}, \ce{Na+(aq)} और \ce{Cl-(aq)}, होते हैं, जो घूमते हैं और
आवेश ढोते हैं। चीनी के \omterm{def:g7:atoms-and-molecules:molecule}{अणुओं} पर कोई आवेश नहीं होता।
\end{solution}

\begin{solution}{exo:g9:ions:5}
\ce{CuSO4}, \ce{AlCl3}, \ce{Na2CO3}।
\end{solution}

\begin{solution}{exo:g9:ions:6}
आयरन(III), \ce{Fe^{3+}}: \ce{Fe^{3+}(aq) + 3OH-(aq) -> Fe(OH)3(s)}।
\end{solution}

\begin{solution}{exo:g9:ions:7}
सल्फ़ेट: 2 अतिरिक्त \omterm{def:g9:inside-the-atom:nucleus}{इलेक्ट्रॉन}। नाइट्रेट: 1।
\end{solution}

\begin{solution}{exo:g9:ions:8}
सिल्वर नाइट्रेट के साथ सफ़ेद: \ce{Cl-}। सोडियम हाइड्रॉक्साइड के साथ नीला:
\ce{Cu^{2+}}। कॉपर(II) क्लोराइड, \ce{CuCl2}।
\end{solution}

\begin{solution}{exo:g9:ions:9}
चार: हर ओर एक (बाएँ, दाएँ, ऊपर, नीचे)।
\end{solution}

\begin{solution}{exo:g9:ions:10}
\ce{Fe^{2+}(aq) + 2OH-(aq) -> Fe(OH)2(s)}।
\end{solution}

\begin{solution}{exo:g9:ions:11}
ज़िंक क्लोराइड: सोडियम \omterm{def:g9:ions:ion}{आयन} सोडियम हाइड्रॉक्साइड के साथ कोई \omterm{def:g9:ions:precipitate}{अवक्षेप} नहीं देते,
ज़िंक \omterm{def:g9:ions:ion}{आयन} सफ़ेद \omterm{def:g9:ions:precipitate}{अवक्षेप} देते हैं:
\ce{Zn^{2+}(aq) + 2OH-(aq) -> Zn(OH)2(s)}।
\end{solution}

\begin{solution}{exo:g9:ions:12}
\ce{CaCl2}: 2000 क्लोराइड \omterm{def:g9:ions:ion}{आयन}। आवेश: $1000 \times (+2) + 2000 \times
(-1) = 0$।
\end{solution}

\begin{solution}{pb:g9:ions:1}
\textbf{1.} आयरन(II), \ce{Fe^{2+}}: \ce{Fe(OH)2}।

\textbf{2.} आयरन(III), \ce{Fe^{3+}}: \ce{Fe(OH)3}।

\textbf{3.} \ce{Fe^{2+}} में एक ज़्यादा है: $26 - 2 = 24$ \omterm{def:g9:inside-the-atom:nucleus}{इलेक्ट्रॉन}, जबकि
\ce{Fe^{3+}} में $26 - 3 = 23$।

\textbf{4.} \ce{Fe^{3+}(aq) + 3OH-(aq) -> Fe(OH)3(s)}।

\textbf{5.} निकलते समय उसमें आयरन(II) \omterm{def:g9:ions:ion}{आयन} होते हैं; \omterm{def:g4:air-a-mixture-of-gases:air}{वायु} के संपर्क में वे
आयरन(III) \omterm{def:g9:ions:ion}{आयन} बन जाते हैं, जो नारंगी ठोस बनाते हैं।

\textbf{6.} \ce{Fe(OH)3}: $(+3) + 3 \times (-1) = 0$।

\textbf{7.} निकलते समय नहीं: लोहे के \omterm{def:g9:ions:ion}{आयन} घुले हुए हैं और छन्ने से निकल जाते
हैं। नारंगी ठोस बन जाने पर हाँ: वह एक \omterm{def:g9:ions:precipitate}{अवक्षेप} है, जिसे छन्ना
रोक लेता है।

\textbf{8.} नहीं: वह एक \omterm{def:g6:pure-substances-and-mixtures:mixture}{मिश्रण} है (पानी और घुले हुए \omterm{def:g9:ions:ion}{आयन})। ताज़ा निकाला गया
और स्वच्छ होने पर वह समांगी है।

\textbf{9.} $56 \times 1.67 \times 10^{-27} = \qty{9.35e-26}{kg}$।

\textbf{10.} $\qty{0.30}{mg} = \qty{0.30e-6}{kg} = \qty{3.0e-7}{kg}$।

\textbf{11.} \omterm{def:g9:ions:ion}{आयन} \omterm{def:g7:atoms-and-molecules:atom}{परमाणु} से दो या तीन \omterm{def:g9:inside-the-atom:nucleus}{इलेक्ट्रॉनों} का अंतर रखता है, जिनका
द्रव्यमान \omterm{def:g9:inside-the-atom:nucleus}{नाभिक} की तुलना में नगण्य है।

\textbf{12.} $3.0 \times 10^{-7} / 9.35 \times 10^{-26} \approx 3.2
\times 10^{18}$ लोहे के \omterm{def:g9:ions:ion}{आयन} प्रति लीटर।
\end{solution}
